\lstset{style=sql}

\section*{Lernziele: SQL}
\begin{todolist}
\item Ich kann zwischen Bits und Bytes unterscheiden und kenne die Bedeutungen von gängigen Grösseneinheiten (Kilo-, Mega-, Giga-, Tera- und Petabyte).
\item Ich weiss betreffend eine Tabelle, was eine Attribut und was eine Zeile ist.
\item Ich kann einfache SQL-Abfragen formulieren mit Befehlen wie \lstinline|SELECT|, \lstinline|WHERE|,  \lstinline|ORDER BY|, \lstinline|LIMIT| und weiteren Befehlen (s. Cheatsheet für vollständige Liste aller zu lernenden SQL-Befehle).
\item Ich kann Aggregatsfunktionen wie \lstinline|SUM|, \lstinline|MIN|, \lstinline|MAX|, \lstinline|COUNT|oder \lstinline|LENGTH| auf einzelne Kolonnen anwenden und kenne deren Bedeutung.
\item Ich kann die oben erwähnten SQL-Begriffe gruppenweise anwenden, indem ich sie mit dem Befehl \lstinline|GROUP BY| kombiniere.
\item Ich kann in eigenen Worten erklären, welche Daten benötigt werden, um ein einfaches soziales Netzwerk zu betreiben und wie diese Daten untereinander zusammenhängen.
\item Ich kann den Unterschied zwischen einem Primär- und einem Frendschlüssel benennen und erklären.
\item Ich kann anhand einer Übersicht von Tabellen in einer Datenbank erkennen, welche Attribute Primär- und welche Sekundärschlüssel sind.
\item Ich kann Informationen aus mehreren Tabellen kombinieren, indem ich den \lstinline|JOIN|-Befehl verwenden.
\item Ich kann Tabellenverbindungen sowohl mit dem Befehl \lstinline|ON| wie \lstinline|USING| erstellen und verstehe, welcher Begriff in welchen Situationen besser geeignet ist.
\item Ich kann aus kombinierten Tabellen neue Informationen gewinnen, indem ich die bereits erwähnten SQL-Codewörter wie z.B. \lstinline|GROUP BY| in Kombination mit einem einem \lstinline|JOIN| verwende.
\item Ich kann Tabellen erstellen, löschen und verändern, indem ich (unter anderem) die Befehle \lstinline|CREATE TABLE|, \lstinline|DROP TABLE| und \lstinline|ALTER TABLE| verwende.
\item Ich kann Einträge (Zeilen) in einer existierenden Tabelle hinzufügen, verändern oder löschen, indem ich die Befehle \lstinline|INSERT INTO|, \lstinline|UPDATE| und \lstinline|DELETE FROM| verwende.
\end{todolist}